\epigraph{Nature speaks in ratios, not in kilograms.}{}

\section{A Historic Moment}


\begin{figure}[ht]
\centering
\begin{tikzpicture}[
  oldbox/.style={draw, rounded corners=3pt, fill=red!8, minimum width=3.2cm, minimum height=1.2cm, text width=3cm, align=center, font=\small},
  newbox/.style={draw, rounded corners=3pt, fill=green!8, minimum width=3.2cm, minimum height=1.2cm, text width=3.2cm, align=center, font=\small},
  arr/.style={-{Stealth[length=3mm]}, thick, blue!60!black},
  lbl/.style={font=\footnotesize\itshape, text=blue!60!black}
]
\node[oldbox] (old) at (-3.5,0) {Prototype kg\\ {\footnotesize Platinum-Iridium}\\ {\footnotesize S\`evres, Paris}};
\node[newbox] (new) at (3.5,0) {$\hbar = 1.054\ldots$\\ $\times 10^{-34}$~J\,s\\ {\footnotesize universal, exact}};
\draw[arr] (old) -- node[above, lbl] {SI Reform 2019} (new);
\node[font=\footnotesize, text=red!60!black] at (-3.5,-1.4) {ages, loses atoms};
\node[font=\footnotesize, text=green!50!black] at (3.5,-1.4) {timeless, universally valid};
\end{tikzpicture}
\caption{The paradigm shift of the SI reform: from physical artefact to universal natural constant.}
\label{fig:si_reform}
\end{figure}
On 20 May 2019 something happened that most people did not notice --- and yet it changed the way we measure the physical world. On that day the fundamental redefinition of the SI base units came into force. Four natural constants were fixed to \emph{exact} values:

\begin{align}
c &= 299\,792\,458 \;\text{m/s} \label{eq:c}\\
\hbar &= 1.054\,571\,817\ldots \times 10^{-34} \;\text{J\,s} \label{eq:hbar}\\
e &= 1.602\,176\,634 \times 10^{-19} \;\text{C} \label{eq:e}\\
k_B &= 1.380\,649 \times 10^{-23} \;\text{J/K} \label{eq:kb}
\end{align}

What does this mean? Until that day the kilogram was defined by a cylinder of platinum-iridium kept in a triply secured vault in Sèvres near Paris. The prototype kilogram --- a physical object that collects dust, loses atoms, ages. The unit of mass was literally a piece of metal.

After the reform the kilogram is defined by the Planck constant~$\hbar$. No longer by an artefact but by a property of the universe itself. Nature became the standard.

\section{The Philosophical Significance}

This redefinition is more than a metrological formality. It carries a deep philosophical message: \emph{The fundamental constants of nature are not measured quantities --- they are definitions.}

If $c$ is exactly~$299\,792\,458$~m/s, then the speed of light defines the metre, not the other way around. If $\hbar$ is exactly fixed, then the quantum of action defines the kilogram. The units are no longer arbitrary conventions --- they are mirrors of nature.

But precisely here a question arises that the SI reform did not answer: \emph{If we take nature as our standard --- why not do so consistently?}

The reform fixed four constants. It left others --- the gravitational constant~$G$, the fine-structure constant~$\alpha$ --- as measured values. There are still constants that we cannot derive from others but must measure. The simplification was a step in the right direction, but it stopped halfway.

\begin{kerngedanke}
The SI reform of 2019 set the direction: away from artefacts, towards natural constants. But it did not dare the logical next step --- namely to ask whether \emph{all} constants could follow from a single one.
\end{kerngedanke}

\section{Ratios Move to Centre Stage}

What the reform implicitly did is something remarkable: it placed \emph{ratios} between natural constants at the centre. When both $c$ and $\hbar$ are exactly fixed, the ratio $c/\hbar$ is likewise exact. And this ratio has physical meaning --- it connects length, time, and energy in a fundamental way.

The question that follows is deceptively simple: if ratios are fundamental --- what happens when we set \emph{all} fundamental constants to one and work solely with ratios?

This question leads us to the next chapter~(\ref{ch:alpha}) --- and to the core of what the T0 theory has to say.
